\SourceLeaf{365}%
Fällen unmittelbarer und kürzer zum Ziele führt, als die Analyse, besonders wenn diese nicht mit Gewandtheit gehandhabt wird, dass jene dann eine ihr eigenthümliche Eleganz hat, wird sie auch besonders in formeller Hinsicht und beim frühern jugendlichen Studium unentbehrlich bleiben, um Einseitigkeit zu verhüten, den Sinn für Strenge und Klarheit zu schärfen, und den Einsichten eine Lebendigkeit und Unmittelbarkeit zu geben, welche durch die analytischen Methoden weit weniger befördert, mitunter eher gefährdet werden. Aus diesen Gründen sieht man mit Vergnügen, dass einige Französische Geometer in den letzten Jahrzehnten angefangen haben, den Theil der Geometrie, welcher sich mit den Verhältnissen von Punkten und Linien, die nicht in Einer Ebene liegen, von verschiedenen Ebenen gegen einander, mit Linien von doppelter Krümmung und mit krummen Flächen beschäftigt, mit besonderer Sorgfalt, und, in so fern dabei bloss geometrische Methoden angewandt werden, als eine besondere Disciplin unter dem Namen der \textit{Géométrie descriptive} zu cultiviren. Dem vorliegenden Werke über diese Wissenschaft müssen wir insbesondere das Lob einer grossen Klarheit und Concision im Vortrage, eines wohlgeordneten Überganges vom Leichtern zum Schwerern, und der Reichhaltigkeit an neuen Ansichten und gelungenen Ausführungen beilegen, und daher das Studium desselben als eine kräftige Geistesnahrung empfehlen, wodurch unstreitig zur Belebung und Erhaltung des echten, in der Mathematik der Neuern sonst manchmal vermissten, geometrischen Geistes viel mit beigetragen werden kann. Ausser dieser rein wissenschaftlichen Seite dieser Untersuchungen kommt auch noch der mannigfaltige Nutzen in Betracht, welchen sie in den Künsten haben, die sich auf Raumverhältnisse beziehen, namentlich in der Zeichenkunst, der Feldmesskunst, der Baukunst, der Befestigungskunst. Auch in dieser Hinsicht hat der Verfasser seine Schrift durch mancherlei Anwendungen interessanter zu machen gewusst, wenn er gleich meistens nur mehr auf sie hingedeutet, als sie wirklich ausgeführt hat.
